%% GENERATED by python/paper/build.py from results/ -- do not edit by hand.
%% Rebuild:  .venv\Scripts\python.exe python\paper\build.py --only US_CRRA_PensChars
%% Source:   results/shocks/US_shocks.csv
\begin{table}[!htb]
\centering
\begin{threeparttable}
\caption{Does CRRA matter for the effect of pension design ($\theta$) in US -- 2020}
\label{table:US:CRRA:pensChars}
\renewcommand{\arraystretch}{1.25}
\begin{tabularx}{.9\textwidth}{YYYYY}
\toprule
\textbf{Scenario} & \textbf{CRRA} ($\rho$) & \textbf{Tax rate} & \textbf{Savings rate} & \textbf{Avg. workweek} \\
\midrule 
Baseline &  & 14.43\% & 21.96\% & 39.39 \\[.5em]\hline\\[-.75em]
 & 0.5 & 19.62\% & 20.47\% & 37.80 \\
$\theta = 0$ & 1.0 & 18.52\% & 20.64\% & 37.98 \\
 & 2.0 & 15.32\% & 21.28\% & 38.44 \\[.5em]\hline\\[-.75em]
 & 0.5 & 11.45\% & 24.44\% & 39.95 \\
$\theta = 1$ & 1.0 & 12.28\% & 22.88\% & 39.86 \\
 & 2.0 & 13.71\% & 21.35\% & 39.73 \\[.5em]\hline\\[-.75em]
\bottomrule
\end{tabularx}
\begin{tablenotes}
\footnotesize
\item The baseline row is at $\rho = 1.0$, the log case, where the calibration targets are hit exactly. Every $\rho$ is separately calibrated (results/calibration/US\_rhoGrid.csv) and its workweek is normalised against its own baseline, so the columns are comparable down the table.
\end{tablenotes}
\end{threeparttable}
\end{table}
